\appendix
\chapter{Phụ lục}

\section{Khai báo sử dụng công cụ AI}

Theo yêu cầu công khai công cụ hỗ trợ, nhóm có sử dụng Claude và OpenAI Codex
để gợi ý cấu trúc, hỗ trợ rà soát mã nguồn, kiểm thử và biên tập diễn đạt.
Nhóm chịu trách nhiệm lựa chọn phương pháp, chạy chương trình, kiểm tra kết quả
với dữ liệu CrisisMMD, đối chiếu metric và xác nhận nội dung báo cáo. Nội dung
do công cụ sinh ra không được xem là bằng chứng nếu không có kết quả thực thi
hoặc nguồn tham khảo tương ứng.

\section{Bảng tự đánh giá}

\begin{table}[H]\centering\small
\caption{Tự đánh giá theo các nhóm tiêu chí của bài tập.}
\begin{tabularx}{\textwidth}{XrrX}
\toprule
\textbf{Tiêu chí} & \textbf{Tối đa} & \textbf{Tự đánh giá} & \textbf{Căn cứ}\\
\midrule
Hiểu bài toán và quyết định & 15 & 14 & Chương 1--2, stakeholder và năm quyết định\\
Dữ liệu và tiền xử lý & 15 & 15 & Nhãn chính thức, feature engineering, leakage audit\\
EDA và insight & 20 & 19 & Train-only, event, K sweep, Apriori, ảnh, conflict\\
Mô hình và DSS & 20 & 18 & Sáu model, tuning, fusion, robustness, policy test\\
Diễn giải và khuyến nghị & 15 & 14 & Chương 6, capacity, giới hạn bằng chứng\\
Báo cáo và minh chứng & 10 & 9 & PDF, notebook có output, dashboard, metrics\\
Tổ chức nhóm & 5 & 2 & Đã có danh sách; phân công/tỷ lệ cần nhóm xác nhận\\
\midrule
\textbf{Tổng} & \textbf{100} & \textbf{91} & Không tự chấm tối đa ở phần chưa hoàn tất\\
\bottomrule
\end{tabularx}
\end{table}

\section{Cấu hình và siêu tham số chính}

\begin{table}[H]\centering\small
\begin{tabularx}{\textwidth}{p{3.5cm}X}
\toprule
\textbf{Thành phần} & \textbf{Cấu hình đã khóa}\\
\midrule
TF--IDF & 2.000 đặc trưng, word n-gram (1,2), stopword tiếng Anh, \(L_2\).\\
Text informative & Linear SVM \(C=3\), class weight balanced, sigmoid calibration 3-fold.\\
Text humanitarian & Logistic Regression \(C=1\), class weight balanced.\\
Image informative & k-NN \(k=41\), distance weighting, Euclidean trên CLIP normalized.\\
Image humanitarian & Logistic Regression \(C=1\), class weight balanced.\\
Decision Tree & max depth 20 trong vòng so sánh thuật toán.\\
Random Forest & 200 cây, balanced subsample.\\
K-Means & \(K=2,\ldots,12\), \(n\_init=10\); \(K=8\) để đối chiếu taxonomy.\\
Apriori & min support .005, min confidence .25 trên transaction đã lọc.\\
CLIP & \code{openai/clip-vit-base-patch32}, frozen, embedding 512 chiều.\\
Late Fusion & Informative text/image .70/.30, threshold .38; category .55/.45.\\
Manual Review & Conflict threshold .54, capacity dev tối đa 25\%.\\
\bottomrule
\end{tabularx}
\end{table}

\section{Tệp minh chứng và khả năng tái lập}

\begin{itemize}[leftmargin=*]
  \item Notebook đã thực thi:
        \code{notebooks/02\_eda.ipynb} và \code{03\_modeling.ipynb}.
  \item Kiểm tra dữ liệu:
        \code{reports/metrics/data\_quality\_summary.json} và
        \code{data/processed/split\_integrity.csv}.
  \item Tuning và kết quả:
        \code{reports/metrics/*\_tuning.csv},
        \code{tuned\_model\_summary.json} và các bảng fusion.
  \item Robustness:
        canonical/robust comparison, bootstrap, event và class stability trong
        \code{reports/metrics/}.
  \item DSS:
        \code{dss\_scenarios.csv},
        \code{dss\_threshold\_sensitivity.csv}.
  \item Kiểm thử: 41 unit test, 6 Streamlit AppTest và integration ảnh thật.
  \item Mã nguồn:
        \url{https://github.com/dafiiscoding/ml-dss}.
\end{itemize}

\section{Thông tin còn cần nhóm xác nhận}

Tên, MSSV, giảng viên và nhóm đã được điền ở đầu báo cáo. Trước khi nộp chính
thức, nhóm cần phân công các gói công việc và điền tỷ lệ đóng góp thực tế tại
\code{docs/team\_info.json}; tổng tỷ lệ của bốn thành viên cần bằng 100\%.
